% !TeX document-id = {0c096819-60cf-4d9a-a18e-64c3d748af84}
% !TeX TXS-program:compile = txs:///pdflatex/[-shell-escape]

\documentclass[11pt]{beamer}
\usepackage[utf8]{inputenc}
\usepackage[T1,T2A]{fontenc}
\usepackage[russian]{babel}
\usepackage{color}
\usepackage{listings}
\usepackage{calc}
\usepackage{graphicx}
\usepackage{epstopdf}
\usepackage{hyperref}
\hypersetup{unicode,colorlinks}
\usepackage{csquotes}
\usepackage{upquote}
\usepackage{cprotect}
\usetheme[progressbar=head,numbering=fraction,block=fill]{metropolis}
\usepackage{dejavu}
\usepackage{etoolbox}
\usepackage{bm}
\usepackage{minted}
%\usepackage[defaultsans]{droidsans}
%\usepackage[defaultmono]{droidmono}
%\makeatletter
%\input droid/t1fdm.fd
%\makeatother

\makeatletter
\DeclareRobustCommand*{\bfseries}{%
    \not@math@alphabet\bfseries\mathbf
    \fontseries\bfdefault\selectfont
    \boldmath
}
\makeatother

\makeatletter
\chardef\straightquote@code=\catcode`'
\chardef\backquote@code=\catcode``
\catcode`'=\active \catcode``=\active
\patchcmd{\@noligs}
{\textasciigrave}
{\fixedtextasciigrave}
{}{}
\newcommand{\fixedtextasciigrave}{%
    \makebox[.5em]{\fontencoding{TS1}\fontfamily{fvs}\selectfont\textasciigrave}% Vera Sans
}
\catcode\lq\'=\straightquote@code
\catcode\lq\`=\backquote@code
\makeatletter

\lstset{
    language=Haskell,
    inputencoding=utf8,
    extendedchars=true,
    breaklines=true,
    escapeinside=!!,
    tabsize=4,
    breakatwhitespace=true,
    keepspaces=true,
    basicstyle=\ttfamily,
    commentstyle=\color{gray}\itshape,
    keywordstyle=
    % keywordstyle=\color{blue}, % TODO https://tex.stackexchange.com/questions/415777/avoid-highlighting-keywords-following-certain-words-in-listings
    % addkeywords={Type}
    % deletekeywords={Monad, String, return, array, show},
    % otherkeywords={::, ->}
}
\lstset{
    literate={а}{{\selectfont\char224}}1
    {б}{{\selectfont\char225}}1
    {в}{{\selectfont\char226}}1
    {г}{{\selectfont\char227}}1
    {д}{{\selectfont\char228}}1
    {е}{{\selectfont\char229}}1
    {ё}{{\"e}}1
    {ж}{{\selectfont\char230}}1
    {з}{{\selectfont\char231}}1
    {и}{{\selectfont\char232}}1
    {й}{{\selectfont\char233}}1
    {к}{{\selectfont\char234}}1
    {л}{{\selectfont\char235}}1
    {м}{{\selectfont\char236}}1
    {н}{{\selectfont\char237}}1
    {о}{{\selectfont\char238}}1
    {п}{{\selectfont\char239}}1
    {р}{{\selectfont\char240}}1
    {с}{{\selectfont\char241}}1
    {т}{{\selectfont\char242}}1
    {у}{{\selectfont\char243}}1
    {ф}{{\selectfont\char244}}1
    {х}{{\selectfont\char245}}1
    {ц}{{\selectfont\char246}}1
    {ч}{{\selectfont\char247}}1
    {ш}{{\selectfont\char248}}1
    {щ}{{\selectfont\char249}}1
    {ъ}{{\selectfont\char250}}1
    {ы}{{\selectfont\char251}}1
    {ь}{{\selectfont\char252}}1
    {э}{{\selectfont\char253}}1
    {ю}{{\selectfont\char254}}1
    {я}{{\selectfont\char255}}1
    {А}{{\selectfont\char192}}1
    {Б}{{\selectfont\char193}}1
    {В}{{\selectfont\char194}}1
    {Г}{{\selectfont\char195}}1
    {Д}{{\selectfont\char196}}1
    {Е}{{\selectfont\char197}}1
    {Ё}{{\"E}}1
    {Ж}{{\selectfont\char198}}1
    {З}{{\selectfont\char199}}1
    {И}{{\selectfont\char200}}1
    {Й}{{\selectfont\char201}}1
    {К}{{\selectfont\char202}}1
    {Л}{{\selectfont\char203}}1
    {М}{{\selectfont\char204}}1
    {Н}{{\selectfont\char205}}1
    {О}{{\selectfont\char206}}1
    {П}{{\selectfont\char207}}1
    {Р}{{\selectfont\char208}}1
    {С}{{\selectfont\char209}}1
    {Т}{{\selectfont\char210}}1
    {У}{{\selectfont\char211}}1
    {Ф}{{\selectfont\char212}}1
    {Х}{{\selectfont\char213}}1
    {Ц}{{\selectfont\char214}}1
    {Ч}{{\selectfont\char215}}1
    {Ш}{{\selectfont\char216}}1
    {Щ}{{\selectfont\char217}}1
    {Ъ}{{\selectfont\char218}}1
    {Ы}{{\selectfont\char219}}1
    {Ь}{{\selectfont\char220}}1
    {Э}{{\selectfont\char221}}1
    {Ю}{{\selectfont\char222}}1
    {Я}{{\selectfont\char223}}1
    {`}{{{\fixedtextasciigrave}}}1
}

% \usemintedstyle{solarized-light}
\newminted[haskell]{haskell}{
    escapeinside=``,
    beameroverlays=true,
    texcomments=true,
    autogobble=true,
    fontsize=\small
}
\newmintinline[haskinline]{haskell}{
    escapeinside=``,
    beameroverlays=true,
    texcomments=true
}

\author{Алексей Романов}
\subtitle{Функциональное программирование на Haskell}
%\logo{}
\institute{МИЭТ}
\subject{Функциональное программирование на Haskell}
%\setbeamercovered{transparent}
%\setbeamertemplate{navigation symbols}{}


\title{Лекция 10: монады}

\begin{document}
\begin{frame}[plain]
  \maketitle
\end{frame}

\begin{frame}[fragile]
  \frametitle{Монады}
  \begin{itemize}
    \item Монады расширяют возможности аппликативных функторов так же, как те расширяют возможности функторов.
    \item А именно, там добавляется
          \begin{haskell}
            class Applicative m => Monad m where
              (>>=) :: m a -> (a -> m b) -> m b
              (>>) :: m a -> m b -> m b
              mx >> my =!\pause! mx >>=!\pause! \_ -> my
          \end{haskell}
          \pause
    \item Для законных экземпляров \haskinline|mx >> my| $=$ \haskinline|mx *> my|.
    \item Историческое имя \haskinline|pure| "--- \haskinline|return|.
          \haskinline|return| не завершает функцию или \haskinline|do|-блок.
  \end{itemize}
\end{frame}

\begin{frame}[fragile]
  \frametitle{От функторов к монадам}
  \begin{itemize}
    \item Больше похоже на знакомые типы, если поменять аргументы \haskinline|(>>=)| местами:
          \begin{haskellsmall}
            fmap  ::   (a -> b) -> f a -> f b
            (<*>) :: f (a -> b) -> f a -> f b
            (=<<) :: (a -> m b) -> m a -> m b
          \end{haskellsmall}
  \end{itemize}
\end{frame}

\begin{frame}[fragile]
  \frametitle{Что можно сделать с монадами}
  \begin{itemize}
    \item В прошлой лекции:
          \emph{\enquote{Структура результата аппликативного вычисления зависит только от структуры аргументов, а не от значений внутри них.}}
    \item Монады снимают это ограничение. Теперь можно написать
          \begin{haskellsmall}
            ifM :: Monad m => m Bool -> m a -> m a -> m a
            ifM mCond mThen mElse =
              mCond >>=!\pause! \cond ->!\pause!
                if cond then mThen else mElse
          \end{haskellsmall}
          \begin{haskellsmall}
            ghci> ifM [True] [1] [2,3] !\pause!
            [1]
            ghci> ifM [False] [1] [2,3] !\pause!
            [2,3]
            ghci> ifM [True, False] [1] [2,3] !\pause!
            [1,2,3]
          \end{haskellsmall}
  \end{itemize}
\end{frame}


\begin{frame}[fragile]
  \frametitle{Стрелки Клейсли}
  \begin{itemize}
    \item Тип второго аргумента \haskinline|(>>=)|, \haskinline|a -> m b|, где \haskinline|m| "--- монада, называется \emph{стрелкой Клейсли}.
    \item \haskinline|a| и \haskinline|b| могут быть как переменными типа, так и конкретными типами.
    \item Стрелки Клейсли можно композировать (импорт \haskinline|Control.Monad|):
          \begin{haskellsmall}
            (>=>) :: Monad m => (a -> m b) -> (b -> m c) ->
                                (a -> m c)
            f >=> g = \x ->!\pause! f x >>=!\pause! g
          \end{haskellsmall}
    \item Есть и аналоги \haskinline|filter|, \haskinline|map| и т.д., работающие со стрелками Клейсли. Их можно найти в \href{https://hackage.haskell.org/package/base/docs/Control-Monad.html}{документации \haskinline{Control.Monad}}.
  \end{itemize}
\end{frame}

\begin{frame}[fragile]
  \frametitle{Законы монад}
  \begin{itemize}
    \item Левая единица:
          \begin{haskell}
            pure x >>= f !$=$! f x
          \end{haskell}
    \item Правая единица:
          \begin{haskell}
            mx >>= pure !$=$! mx
          \end{haskell}
    \item Ассоциативность:
          \begin{haskellsmall}
            (mx >>= f) >>= g !$=$! mx >>= (\x -> f x >>= g)
          \end{haskellsmall}
    \item Кроме этого, должно быть согласование с \haskinline|<*>|:
          \begin{haskellsmall}
            mf <*> mx !$=$!
              mf >>= (\f -> mx >>= (\x -> pure (f x)))
          \end{haskellsmall}
          Удобная запись немного позже.
    \item Может быть, они не совсем интуитивны.
  \end{itemize}
\end{frame}

\begin{frame}[fragile]
  \frametitle{Законы монад: композиция}
  \begin{itemize}
    \item С помощью монадической композиции они записываются естественнее:
    \item \haskinline|pure >=> f| $=$ \haskinline|f| \quad (левая единица)
    \item \haskinline|f >=> pure| $=$ \haskinline|f| \quad (правая единица)
    \item \haskinline|(f >=> g) >=> h| $=$ \haskinline|f >=> (g >=> h)| \\
          (ассоциативность)
  \end{itemize}
\end{frame}

\begin{frame}[fragile]
  \frametitle{Примеры монад}
  \begin{itemize}
    \item \haskinline|Maybe| "--- монада. Определим:
          \begin{haskell*}{texcomments=false}
            instance Monad Maybe where
              -- (>>=) :: Maybe a -> (a -> Maybe b) ->
              --          Maybe b
              Nothing >>= _ =!\pause! Nothing
              Just x  >>= f =!\pause! f x
          \end{haskell*}
    \item Списки тоже монада:
          \begin{haskell}
            -- (>{}>=) ::\pause{} [a] -> (a -> [b]) -> [b]
            []     >>= _ =!\pause! []
            (x:xs) >>= f =!\pause! f x ++ (xs >>= f)
          \end{haskell}
          Или проще:
          \begin{haskell}
            xs >>= f = [y |!\pause! x <- xs, y <- f x]
          \end{haskell}
  \end{itemize}
\end{frame}

\begin{frame}[fragile]
  \frametitle{Ещё примеры монад}
  \begin{itemize}
    \item \haskinline|instance Monad Identity|, где \haskinline|newtype Identity a = Identity a|.
    \item \haskinline|instance Monad (Either c)|.
    \item \haskinline|instance Monad ((->) c)|: функции с фиксированным типом аргумента.
    \item \haskinline|instance Monoid c => Monad ((,) c)|: пары с фиксированным типом первого элемента, если этот тип "--- моноид.
  \end{itemize}
\end{frame}

\begin{frame}[fragile]
  \frametitle{\ldots и не монад}
  \begin{itemize}
    \item Есть и примеры аппликативных функторов, для которых нельзя определить экземпляр монады, удовлетворяющий законам:
    \item
          \begin{haskell}
            newtype ConstInt a = ConstInt Int
            fmap f (ConstInt x) = ConstInt x
            pure _ = ConstInt 0
            ConstInt x <*> ConstInt y = ConstInt (x + y)
          \end{haskell}
    \item Легко увидеть, что \haskinline|pure x >>= f| $=$ \haskinline|f x| не может выполняться ни для какого определения \haskinline|>>=|: правая часть зависит от \haskinline|x|, а левая нет.
          \pause
    \item Ещё пример "--- \haskinline|ZipList| (если не допускать \emph{только} бесконечные списки, как в домашнем задании, или активное использование $\bot$).
  \end{itemize}
\end{frame}

\begin{frame}[fragile]
  \frametitle{\haskinline|do|-нотация}
  \begin{itemize}
    \item Для монад есть специальный синтаксис, которым часто удобнее пользоваться.
    \item Скажем, у нас есть цепочка операций
          \begin{haskellsmall}
            action1 >>= (\x1 -> action2 x1 >>=
              (\x2 -> action3 x1 x2 >> action4 x1))
          \end{haskellsmall}
    \item
          Сначала перепишем так:
          \begin{haskellsmall}
            action1 >>= \x1 ->
              action2 x1 >>= \x2 ->
              action3 x1 x2 >>
              action4 x1
          \end{haskellsmall}
    \item
          В \haskinline|do|-блоке строки вида \haskinline|action1 >>= \x1 ->| превращаются в \haskinline|x1 <- action1|, а \haskinline|>>| пропадает:
          \begin{haskellsmall}
            do x1 <- action1
               x2 <- action2 x1
               action3 x1 x2
               action4 x1
          \end{haskellsmall}
  \end{itemize}
\end{frame}

\begin{frame}[fragile]
  \frametitle{Законы монад в \haskinline|do|-нотации}
  \begin{itemize}[<+->]
    \item Законы монад также можно записать через \haskinline|do|:
    \item
          \begin{haskell}
            do y <- pure x    !$\mathtt{=}$!    f x
               f y
          \end{haskell}
    \item
          \begin{haskell}
            do x <- mx        !$\mathtt{=}$!    mx
               pure x
          \end{haskell}
    \item
          \begin{haskell}
            do x <- mx        !$\mathtt{=}$!    do y <- do x <- mx
               y <- f x       !\textcolor{white}{$\mathtt{=}$}!               f x
               g y            !\textcolor{white}{$\mathtt{=}$}!       g y
          \end{haskell}
    \item
          \begin{haskell}
            mf <*> mx         !$\mathtt{=}$!    do f <- mf
                              !\textcolor{white}{$\mathtt{=}$}!       x <- mx
                              !\textcolor{white}{$\mathtt{=}$}!       pure (f x)
          \end{haskell}
  \end{itemize}
\end{frame}

\begin{frame}[fragile]
  \frametitle{Общая форма \haskinline|do|-нотации}
  \begin{itemize}
    \item Каждая строка \haskinline|do|-блока имеет вид \haskinline|образец <- м_выражение|, \haskinline|let образец = выражение| или просто \haskinline|м_выражение|.
    \item Первые два вида не могут быть в конце.
    \item \haskinline|м_выражение| должно иметь тип \haskinline|m a| для какой-то монады \haskinline|m| и типа \haskinline|a|.
    \item \haskinline|m| одна для всех строк, \haskinline|a| могут различаться.
    \item \haskinline|образец| в строке с \haskinline|<-| имеет тип \haskinline|a|.
    \item Если образец после \haskinline|<-| может не совпасть, нужен экземпляр \haskinline|MonadFail|; при несовпадении вызывается \haskinline|fail|.
  \end{itemize}
\end{frame}

\begin{frame}[fragile]
  \frametitle{Несовпадение образца в \haskinline|do|-блоке}
  \begin{itemize}
    \item Для \haskinline|Maybe| несовпадение образца даёт \haskinline|Nothing|:
          \begin{haskell}
            safeHead :: Maybe [Int] -> Maybe Int
            safeHead mx = do
              (x : _) <- mx
              pure x
          \end{haskell}
          \pause
          \begin{haskell}
            ghci> safeHead (Just [2, 3])
            !\pause!Just 2
            ghci> safeHead (Just [])
            !\pause!Nothing
          \end{haskell}
    \item При несовпадении \haskinline|pure x| не выполняется.
    \item Подробности "--- в \href{https://hackage.haskell.org/package/base/docs/Control-Monad-Fail.html#t:MonadFail}{документации \haskinline{MonadFail}}.
    \item Дальше нам достаточно образцов, которые всегда совпадают.
  \end{itemize}
\end{frame}

\begin{frame}[fragile]
  \frametitle{Функции над произвольными монадами}
  \begin{itemize}[<+->]
    \item В \haskinline|Prelude| есть \haskinline|(=<<)|, а также:
          \begin{haskellsmall}
            sequence :: (Traversable t, Monad m) =>
              t (m a) -> m (t a)

            mapM :: (Traversable t, Monad m) =>
              (a -> m b) -> t a -> m (t b)
          \end{haskellsmall}
    \item \haskinline|sequenceA| и \haskinline|traverse| из прошлой лекции делают то же самое, но требуют только \haskinline|Applicative m|.
    \item Подставьте \haskinline|t = []|, а затем конкретные монады (например, \haskinline|Maybe| и \haskinline|[]|). Что будут делать эти функции?
  \end{itemize}
\end{frame}

\begin{frame}[fragile]
  \frametitle{Функции из \haskinline|Control.Monad|}
  \begin{itemize}[<+->]
    \item Нужен импорт \haskinline|Control.Monad|: здесь определены \haskinline|(>=>)|, \haskinline|join| и \haskinline|zipWithM|.
    \item
          \begin{haskell}
            join :: Monad m => m (m a) -> m a
          \end{haskell}
          Эту функцию можно было бы взять как базовую и выразить \haskinline|>>=| через неё.
    \item
          \begin{haskell}
            zipWithM :: Applicative m =>
              (a -> b -> m c) -> [a] -> [b] -> m [c]
          \end{haskell}
    \item Подумайте, что они делают для \haskinline|Maybe| и \haskinline|[]|.
  \end{itemize}
\end{frame}

\begin{frame}[fragile]
  \frametitle{\haskinline|join| и \haskinline|>>=|: упражнение}
  \begin{itemize}
    \item Выразите \haskinline|join| через \haskinline|>>=|, а затем \haskinline|>>=| через \haskinline|fmap| и \haskinline|join|:
          \begin{haskell}
            join :: Monad m => m (m a) -> m a
            join mm = ???

            mx >>= g = ???
          \end{haskell}
    \item Подсказка: подумайте о типе \haskinline|fmap g mx|.
  \end{itemize}
\end{frame}

\begin{frame}[fragile]
  \frametitle{Монада \haskinline|State|}
  \begin{itemize}
    \item Рассмотрим ещё один, более сложный пример:
    \item
          \begin{haskellsmall}
            newtype State s a = State { runState :: s -> (a, s) }
          \end{haskellsmall}
          Это \enquote{вычисления с состоянием}, которые могут выдать результат, зависящий от состояния, и изменить это состояние.
    \item \haskinline|runState| принимает вычисление и начальное состояние, возвращает результат и конечное состояние.
  \end{itemize}
\end{frame}

\begin{frame}[fragile]
  \frametitle{\haskinline|State|: экземпляры Functor и Applicative}
  \begin{itemize}
    \item Как сделать эти вычисления аппликативным функтором?
          \begin{haskellsmall}
            instance Functor (State s) where
              fmap f (State gx) = State (\s1 ->!\pause!
                let (x, s2) = gx s1
                in!\pause! (f x, s2))

            instance Applicative (State s) where
              pure x = State (\s1 ->!\pause! (x, s1))
              (State gf) <*> (State gx) = State (\s1 ->
                let (f, s2) = gf s1!\pause!
                    (x, s3) = gx s2
                in!\pause! (f x, s3))
          \end{haskellsmall}
  \end{itemize}
\end{frame}

\begin{frame}[fragile]
  \frametitle{Монада \haskinline|State|}
  \begin{itemize}
    \item А теперь монадой:
          \begin{haskellsmall}
            instance Monad (State s) where
              State gx >>= f = State (\s1 ->
                let (x, s2) = gx s1!\pause!
                    State gy = f x
                in!\pause! gy s2)
          \end{haskellsmall}
          \pause
    \item Определим вспомогательные функции:
          \begin{haskellsmall}
            get :: State s s
            get = State!\pause! (\s -> (s, s))

            put :: s -> State s ()
            put x = State!\pause! (\_ -> ((), x))

            modify :: (s -> s) -> State s ()
            modify f = do x <- get
                          put (f x)
          \end{haskellsmall}
  \end{itemize}
\end{frame}

\begin{frame}[fragile]
  \frametitle{\haskinline|State|: пример вычисления}
  \begin{itemize}
    \item Прочитаем состояние, увеличим его и вернём оба значения:
          \begin{haskell}
            increment :: State Int (Int, Int)
            increment = do
              old <- get
              modify (+1)
              new <- get
              pure (old, new)
          \end{haskell}
          \pause
    \item Запустим вычисление с начальным состоянием \haskinline|10|:
          \begin{ghci}
            ghci> runState increment 10
            ((10,11),11)
          \end{ghci}
    \item \haskinline|(10,11)| "--- результат, последнее \haskinline|11| "--- конечное состояние.
  \end{itemize}
\end{frame}

\begin{frame}[fragile]
  \frametitle{\haskinline|State| и \haskinline|traverse|: упражнение}
  \begin{itemize}
    \item Дополнительное упражнение: пронумеруйте элементы списка с помощью состояния.
    \item Каждый элемент получает текущий номер, затем счётчик увеличивается:
          \begin{haskell}
            number :: a -> State Int (Int, a)
            number x = ???
          \end{haskell}
    \item Подсказка: используйте \haskinline|get|, \haskinline|modify| и \haskinline|pure|.
          \pause
    \item Ожидаемый результат обхода:
          \begin{ghcismall}
            ghci> runState (traverse number ["a", "b"]) 0
            ([(0,"a"),(1,"b")],2)
          \end{ghcismall}
  \end{itemize}
\end{frame}

\begin{frame}[fragile]
  \frametitle{Монада \haskinline|State| в библиотеках}
  \begin{itemize}
    \item \hackage{transformers} содержит сами трансформеры монад (\haskinline|StateT|, \haskinline|ReaderT|, \ldots; о них позже в курсе), модуль \href{https://hackage.haskell.org/package/transformers/docs/Control-Monad-Trans-State.html}{\haskinline{Control.Monad.Trans.State}}.
    \item \hackage{mtl} строится на \hackage{transformers} и добавляет классы вроде \href{https://hackage.haskell.org/package/mtl/docs/Control-Monad-State-Class.html#t:MonadState}{\haskinline{MonadState}}, чтобы \haskinline|get| и \haskinline|put| работали без явного \haskinline|lift|; модуль \href{https://hackage.haskell.org/package/mtl/docs/Control-Monad-State.html}{\haskinline{Control.Monad.State}}.
    \item В обоих \haskinline|State s| определён как \haskinline|StateT s Identity|.
    \item Поэтому \emph{конструктора} \haskinline|State| там нет: вместо него функция
          \begin{haskell}
            state :: (s -> (a, s)) -> State s a
          \end{haskell}
    \item \haskinline|runState|, \haskinline|get|, \haskinline|put|, \haskinline|modify| есть, а также \haskinline|evalState| (только результат) и \haskinline|execState| (только состояние).
  \end{itemize}
\end{frame}

\begin{frame}[fragile]
  \frametitle{Дополнительное чтение}
  \begin{itemize}
    \item \href{https://www.snoyman.com/blog/2017/01/functors-applicatives-and-monads}{Functors, Applicatives, and Monads}
    \item \href{https://wiki.haskell.org/Typeclassopedia}{Typeclassopedia} (ещё раз)
    \item \href{https://smunix.github.io/dev.stephendiehl.com/hask/index.html#monads}{What I Wish I Knew When Learning Haskell: Monads}
    \item \href{https://wiki.haskell.org/All_About_Monads}{All About Monads}
    \item И более сложные монады:
          \begin{itemize}
            \item \href{http://blog.sigfpe.com/2008/12/mother-of-all-monads.html}{The Mother of all Monads}
            \item \href{https://lukepalmer.wordpress.com/2008/08/10/mindfuck-the-reverse-state-monad/}{Mindfuck: The Reverse State Monad}
            \item \href{https://web.archive.org/web/2024/https://tech-blog.capital-match.com/posts/5-the-reverse-state-monad.html}{The Curious Time-Traveling Reverse State Monad}
          \end{itemize}
  \end{itemize}
\end{frame}

\end{document}
